\documentclass[11pt]{article}
\usepackage[margin=1in]{geometry}
\usepackage{amsmath,amssymb,amsthm}
\usepackage{microtype}

\title{Derivation of the two-direction superellipse from the response form}
\date{}

\begin{document}
\maketitle

\noindent Starting point:
\[
R(v) \;=\; F\!\left(\sum_i w_i\, |d_i^\top v|^n\right).
\]

\medskip
\noindent\textbf{Step 1: set up the perturbation.} Perturb along the unit direction
\[
\hat v(\varphi) \;=\; \cos\varphi\, d_1 + \sin\varphi\, d_2,
\]
assuming $d_1 \perp d_2$ (so $\|\hat v(\varphi)\|=1$). The perturbation angle $\theta$ parameterises the norm-matched displacement on the sphere, so the perturbation is $\sin\theta\cdot\hat v(\varphi)$. The projections onto the two directions are
\[
d_1^\top\bigl(\sin\theta\,\hat v\bigr) = \sin\theta\cos\varphi, \qquad d_2^\top\bigl(\sin\theta\,\hat v\bigr) = \sin\theta\sin\varphi.
\]
For the small angles relevant to breaking, $\sin\theta \approx \theta$, so we use $\theta\cos\varphi$ and $\theta\sin\varphi$ directly.

\medskip
\noindent\textbf{Step 2: plug into the response.} Assume the remaining features $d_i$ ($i \neq 1,2$) have negligible overlap with $\hat v(\varphi)$ --- the relevant approximation when $d_1, d_2$ are isolated feature directions. Then the inner sum collapses to two terms:
\[
R \;=\; F\!\left(w_1 (\theta\cos\varphi)^n + w_2 (\theta\sin\varphi)^n\right)
  \;=\; F\!\left(\theta^n\bigl[w_1 \cos^n\varphi + w_2 \sin^n\varphi\bigr]\right).
\]

\medskip
\noindent\textbf{Step 3: breaking condition.} The plateau breaks when $F$ crosses a fixed level $F^*$, equivalently when its argument crosses a fixed scalar $\xi^*$. The breaking angle $\theta(\varphi)$ satisfies
\begin{equation}
\theta(\varphi)^n\bigl[w_1 \cos^n\varphi + w_2 \sin^n\varphi\bigr] \;=\; \xi^*. \label{eq:break}
\end{equation}

\medskip
\noindent\textbf{Step 4: calibrate $w_1, w_2$ from the two axis sweeps.} Along $d_1$ alone ($\varphi = 0$), \eqref{eq:break} gives $\theta_1^n\, w_1 = \xi^*$, hence
\[
w_1 \;=\; \xi^*/\theta_1^n.
\]
Along $d_2$ alone ($\varphi = \pi/2$), similarly
\[
w_2 \;=\; \xi^*/\theta_2^n.
\]

\medskip
\noindent\textbf{Step 5: substitute and simplify.} Plug $w_1, w_2$ into \eqref{eq:break} and divide by $\xi^*$:
\[
\theta(\varphi)^n\left[\frac{\cos^n\varphi}{\theta_1^n} + \frac{\sin^n\varphi}{\theta_2^n}\right] \;=\; 1,
\]
which rearranges to
\[
\boxed{\;\left(\frac{\theta(\varphi)\cos\varphi}{\theta_1}\right)^{\!n} + \left(\frac{\theta(\varphi)\sin\varphi}{\theta_2}\right)^{\!n} \;=\; 1.\;}
\]
This is a superellipse of exponent $n$ in the normalised coordinates
\[
\bigl(\theta(\varphi)\cos\varphi / \theta_1,\ \theta(\varphi)\sin\varphi / \theta_2\bigr).
\]

\medskip
\noindent\textbf{Assumptions used.}
\begin{itemize}
\item $d_1 \perp d_2$, so $\hat v(\varphi)$ is unit norm and the two projections are independent.
\item Other features $d_i$ have negligible overlap with $\hat v(\varphi)$ for all $\varphi$, so they drop out of the inner sum.
\item Small-angle regime, $\sin\theta \approx \theta$, so the sphere parameterisation coincides with the angle used in the final expression.
\item $F$ is strictly monotonic through the threshold region, so ``breaking'' corresponds to a fixed value $\xi^*$ of its argument.
\end{itemize}

\noindent The first two assumptions are what make the two-direction fit work pair-by-pair without needing the full feature catalogue: within a single pair's subspace, only $w_1, w_2$ appear, and both are absorbed into $\theta_1, \theta_2$. The superellipse shape is then purely a statement about $n$.

\end{document}
